
\chapter*{Prakata}

\section*{Tentang Tea Time Numerical Analysis}

\addcontentsline{toc}{chapter}{Prakata}
\addcontentsline{toc}{section}{Tentang Tea Time Numerical Analysis}
\fancyhead[RE]{Prakata}
\fancyhead[LO]{}

Salam! Terima kasih telah meluangkan waktu untuk membaca \textit{Tea
Time Numerical Analysis}. Buku teks ini lahir dari keinginan untuk
menyumbangkan buku teks pengantar Analisis Numerik yang layak digunakan
dan sepenuhnya gratis bagi dosen serta mahasiswa matematika. Ketika
proyek ini dimulai (musim panas 2012), tersedia buku-buku teks terbitan
konvensional (bersampul keras dan sangat mahal), terutama buku yang
sangat baik, \textit{Numerical Analysis} karya Burden dan Faires,
yang saat itu telah mencapai edisi kesembilan. Seperti yang dapat
Anda duga dari banyaknya edisi tersebut, buku ini merupakan karya
klasik. Buku itu termasuk salah satu dari sedikit buku teks analisis
numerik yang ditujukan kepada matematikawan, bukan ilmuwan atau insinyur.
Bahkan, saya belajar dari salah satu edisi awalnya pada pertengahan
tahun 1990-an! Pada musim panas 2012 juga tersedia beberapa situs
web gratis, terutama situs populer \href{http://nm.mathforcollege.com/}{http://nm.mathforcollege.com/},
lengkap dengan kuliah video. Namun, dari semua sumber yang dapat saya
temukan, tidak ada yang menyediakan buku teks lengkap dalam satu berkas
PDF yang dapat diunduh dan dirancang untuk perkuliahan matematika.
Menjadi sumber semacam itulah tujuan utama \textit{Tea Time Numerical
Analysis}.

Ungkapan ``waktu minum teh'' tidak sekadar dimaksudkan untuk memberi
buku ini judul yang menarik. Ungkapan tersebut dimaksudkan untuk menggambarkan
corak umum pembahasan di dalamnya. Sebagian besar materi akan disajikan
seolah-olah sedang diceritakan kepada seorang mahasiswa ketika minum
teh di kampus, tetapi dengan manfaat dari perencanaan yang saksama.
Tidak akan ada kotak biru besar yang menyoroti pokok-pokok utama,
tidak ada rentetan contoh setelah pengantar singkat suatu topik, dan
tidak ada struktur teorema\ldots bukti\ldots teorema\ldots bukti.
Sebaliknya, istilah, definisi, teorema, dan contoh yang diperlukan
akan dirajut ke dalam gaya yang lebih menyerupai percakapan. Saya
berharap perpaduan matematika formal dan informal ini akan lebih mudah
dicerna dan, berani saya katakan, membuat mahasiswa lebih terdorong
untuk membaca dalam format ini. 

Pembaca yang menyukai penyajian yang lebih lazim mungkin tetap akan
merasa bahwa buku teks ini cukup sesuai dengan selera mereka. Pada
akhir setiap percakapan terdapat ringkasan konsep-konsep utama, dan
sejumlah latihan diselesaikan secara lengkap dan terperinci di lampiran.
Karena itu, penyajian yang lebih mendekati bentuk lazim dapat diperoleh
dengan mencari teorema di dalam percakapan, membaca konsep-konsep
utama, lalu langsung menuju latihan yang dilengkapi solusi. Saya berharap
kebanyakan pembaca tidak memilih cara tersebut, tetapi pilihan itu
tetap tersedia. Bagaimanapun, bagi kebanyakan pembaca, latihan yang
dilengkapi solusi merupakan bacaan penting. Belajar melalui contoh
sering kali merupakan cara yang paling efektif. Setelah membaca suatu
bagian, atau setidaknya meninjaunya sekilas, pembaca sangat dianjurkan
untuk langsung menuju pernyataan latihan yang dilengkapi solusi (ditandai
dengan  \hasasolution  atau  \haspartwithsolution), memikirkan kemungkinan
penyelesaiannya, mengerjakannya jika mampu, lalu membuka bagian belakang
buku untuk melihat penyelesaian lengkap. Dengan menempatkan solusi
di lampiran, harapannya pembaca akan lebih terdorong untuk mencoba
menyelesaikan latihan secara mandiri sebelum melihat solusinya.

Cakupan topik dalam \textit{Tea Time Numerical Analysis} cukup lazim.
Buku ini dimulai dengan bab pengantar, dilanjutkan dengan metode pencarian
akar, interpolasi (bagian 1), kalkulus numerik, interpolasi (bagian
2), dan pada edisi kedua ditambahkan satu bab tentang persamaan diferensial.
Lima bab pertama mencakup materi yang di SCSU membentuk perkuliahan
analisis numerik semester pertama. Karena buku ini dimaksudkan untuk
digunakan sebagai unduhan gratis atau buku cetak sesuai permintaan
yang murah, tidak ada upaya untuk menekan jumlah halaman ataupun membatasi
banyaknya diagram dan warna. Bahkan, saya memanfaatkan biaya penyampaian
yang murah itu sebagai kebebasan untuk melakukan hal yang justru sebaliknya.
Saya menambahkan banyak uraian dan diagram yang tidak mutlak diperlukan
untuk mempelajari analisis numerik, tetapi setidaknya berkaitan secara
tidak langsung dan mungkin menarik bagi sebagian pembaca. Sebagian
besar uraian ini disajikan sebagai selingan sehingga mudah dikenali.
Sebagai contoh, teorema Taylor memegang peranan yang begitu sentral
dalam bidang ini sehingga buku ini tidak hanya menyajikan pernyataannya.
Bukti teorema tersebut dan sedikit sejarah ditambahkan sebagai ``kudapan''.
Tentu saja bagian-bagian itu boleh dilewati, tetapi disertakan agar
pembaca memperoleh pemahaman yang lebih utuh mengenai teorema fundamental
dalam analisis numerik ini. Contoh lainnya: sebagai penggemar sistem
dinamik, saya merasa mustahil untuk tidak menyertakan bagian tentang
visualisasi Metode Newton. Gambar-gambar Metode Newton sebagai sistem
dinamik yang kuat sekaligus indah semestinya mengundang kekaguman
dan pertanyaan, meskipun kepentingannya hanya bersinggungan dengan
pokok bahasan utama. Tentu saja ada contoh lain berupa materi yang
agak kurang penting, tetapi setiap bagian tersebut hadir untuk memperdalam
pemahaman atau penghargaan pembaca terhadap bidang ini, sekalipun
materinya tidak mutlak diperlukan dalam pengantar analisis numerik.
Meskipun versi 2.5 tidak menambahkan bagian baru, versi tersebut memuat
banyak koreksi, latihan baru, dan perubahan uraian. Akibatnya, penomoran
halaman dan latihan berubah.

Implementasi metode numerik dalam bentuk kode komputer juga akan dibahas.
Meskipun seseorang dapat mengabaikan bagian dan latihan pemrograman
tetapi tetap memperoleh manfaat dari buku ini, saya sungguh percaya
bahwa penghargaan penuh terhadap isinya tidak dapat dicapai tanpa
``turun tangan'' dan melakukan pemrograman sendiri. Akan sangat
membantu jika pembaca setidaknya pernah sedikit mengenal pemrograman,
entah Python, Java, C, pemrograman web, atau hampir apa saja. Namun,
saya telah berusaha memberikan perincian yang cukup agar bahkan mereka
yang belum pernah menulis satu baris kode pun dapat mengikuti bagian
pembelajaran ini. Sejalan dengan keinginan untuk menghasilkan pengalaman
belajar yang sepenuhnya gratis, GNU Octave dipilih sebagai bahasa
pemrograman untuk buku ini. GNU Octave (singkatnya Octave) tersedia
secara bebas bagi siapa saja! Perangkat lunak ini bebas diunduh dan
digunakan. Kode sumbernya bebas diunduh dan dipelajari. Siapa pun
juga dipersilakan mengubah atau menambahkan kode jika menghendakinya.
Sebagai manfaat tambahan, pengguna MATLAB yang jauh lebih dikenal
tidak perlu dibebani dengan mempelajari bahasa baru. Octave merupakan
klona MATLAB. Sesuai rancangannya, hampir setiap program yang ditulis
dalam MATLAB akan berjalan di Octave tanpa perubahan. Jadi, jika Anda
memiliki akses ke MATLAB dan lebih suka menggunakannya, Anda dapat
melakukannya tanpa khawatir. Saya telah berupaya keras memastikan
bahwa setiap baris kode Octave dalam buku ini dapat dijalankan apa
adanya di MATLAB. Meskipun demikian, ada kemungkinan sebagian kode
tidak berjalan di MATLAB. Kode tersebut hanya diuji di Octave! Jika
Anda menemukan kode yang tidak berjalan di MATLAB, mohon beri tahu
saya. Versi 2.5 memuat penulisan ulang yang cukup besar atas rincian
terkait Octave, mulai dari instalasi hingga penulisan kode, dan bahkan
mencakup informasi tentang komputasi awan.

Saya berharap Anda menikmati pembacaan \textit{Tea Time Numerical
Analysis}. Menulisnya merupakan suatu kesenangan bagi saya. Masukan
selalu diterima dengan senang hati.

\bigskip{}

\hfill{} \hfill{} %
\begin{tabular}{l}
Leon Q. Brin\tabularnewline
brinl1@southernct.edu\tabularnewline
\end{tabular} \hfill{} 

\section*{Cara Memperoleh GNU Octave}

\addcontentsline{toc}{section}{Cara Memperoleh Octave}

Octave dikembangkan oleh Proyek GNU untuk sistem operasi GNU, yang
paling sering dipasangkan dengan kernel Linux. Karena itu, pada dasarnya
Octave merupakan perangkat lunak GNU/Linux. Octave berjalan secara
natif pada mesin GNU/Linux. Pengguna Linux dapat menemukannya di repositori
perangkat lunak distribusi mereka. Demi kemudahan, beberapa cara lain
untuk mengakses Octave diuraikan di bawah ini.

\subsubsection*{Daring (di awan)}

Mungkin cara termudah untuk mulai menggunakan Octave adalah melalui
peramban web. Pendekatan ini menawarkan keunggulan berupa waktu mulai
yang singkat, tanpa instalasi, dan kebutuhan perangkat keras yang
sangat ringan. Lingkungan Octave daring dapat dijalankan pada stasiun
kerja paling bertenaga maupun pada telepon genggam. Saat tulisan ini
dibuat, tersedia sedikitnya dua platform daring gratis.
\begin{itemize}
\item \href{https://octave-online.net/}{https://octave-online.net/} 
\item \href{https://cocalc.com/}{https://cocalc.com/}
\end{itemize}
Octave Online dirancang khusus sebagai antarmuka Octave dan menyediakan
fitur yang berguna untuk pembelajaran di kelas. Ajukan tiket dukungan
untuk menanyakan penyiapan suatu kelas. Kode kelas yang dihasilkan
memungkinkan pengajar melihat ruang kerja mahasiswa dan berkolaborasi
secara waktu nyata setelah mereka mendaftar ke kelas tersebut. Fitur
bucket memungkinkan pengajar memasang kode hanya-baca untuk diimpor
mahasiswa ke akun mereka, sehingga kode untuk praktikum atau tugas
mudah dibagikan.

CoCalc (Collaborative Calculation in the Cloud) menyediakan akses
ke Octave, R, SageMath, \LaTeX , C++, Julia, dan berbagai kernel lainnya.
CoCalc menghadirkan lingkungan yang jauh lebih lengkap, termasuk fitur
pengelolaan kelas seperti pembuatan kelas, pembagian dan pengumpulan
tugas, serta penyuntingan kolaboratif; secara keseluruhan, tersedia
seperangkat alat minimum untuk menyelenggarakan perkuliahan. Layanan
ini menawarkan akses akun gratis maupun berbayar. Meskipun mereka
memperingatkan bahwa akun gratis mungkin berjalan lambat, akun tersebut
mungkin tetap sesuai dengan kebutuhan Anda. Berdasarkan pengalaman
saya, akun gratis memerlukan cukup banyak kesabaran karena masalah
koneksi dan pelambatan tidak terelakkan.

\subsubsection*{Secara lokal (di komputer pribadi Anda)}

Untuk Windows, unduh penginstal Octave di \href{https://www.gnu.org/software/octave/download\#ms-windows}{https://www.gnu.org/software/octave/download\#ms-windows}.

\medskip{}

\noindent Untuk macOS, lihat \href{https://octave-app.org/}{https://octave-app.org/}.
Mungkin metode instalasi yang paling sederhana adalah mengikuti tautan
menuju halaman wiki, lalu menggunakan tautan di bawah macOS App Bundles.
Jika Anda memakai metode ini, saat pertama kali menjalankan Octave
Anda perlu mengeklik ikon aplikasi sambil menahan tombol Control,
lalu memilih Open pada menu sembul. Jika Anda hanya mengekliknya seperti
aplikasi lain, macOS akan menolak menjalankannya karena aplikasi tersebut
tidak ditandatangani sehingga dianggap tidak tepercaya.

\medskip{}

\noindent Untuk \href{https://www.gnu.org/gnu/linux-and-gnu.en.html}{GNU/Linux},
distribusi Anda kemungkinan memiliki paket bernama octave yang dapat
dipasang dengan cara biasa. Jika paket tersebut tidak dapat ditemukan,
lihat catatan di \href{https://www.gnu.org/software/octave/download\#linux}{https://www.gnu.org/software/octave/download\#linux}.

\medskip{}

\noindent Untuk perincian lebih lanjut, kunjungi situs web GNU Octave.
\begin{center}
\href{https://www.gnu.org/software/octave/}{https://www.gnu.org/software/octave/}
\par\end{center}

\section*{Cara Memperoleh Kode}

\addcontentsline{toc}{section}{Cara Memperoleh Kode}

Seluruh kode yang muncul dalam buku teks ini dapat diunduh dari situs
pendamping buku,

\begin{center}
\url{http://lqbrin.github.io/tea-time-numerical/ancillaries.html}.
\par\end{center}

\noindent Kode yang dicetak di dalam Tea Time Numerical Analysis dan
kode elektronik yang menyertainya didistribusikan berdasarkan GNU
General Public License (GPL). Perinciannya tersedia di situs web tersebut.

\section*{Ucapan Terima Kasih}

\addcontentsline{toc}{section}{Ucapan Terima Kasih}

Dengan penuh rasa terima kasih, saya mengakui dukungan besar yang
saya terima selama penulisan buku teks ini: mulai dari kesabaran keluarga
inti saya---Amy, Cecelia, dan Victorija---ketika perhatian saya
terserap oleh layar laptop, hingga kesediaan para peserta kelas Seminar
Musim Semi 2013, Elizabeth Field, Rachael Ivison, Amanda Reyher, dan
Steven Warner, untuk membaca serta mengkritik versi awal bab pertama.
Staf Woodbridge Public Library, khususnya Pamela Wilonski, turut menyediakan
lingkungan yang tenang dan inspiratif untuk menulis sebagian besar
teks edisi pertama. Terima kasih sebesar-besarnya kepada Dick Pelosi
atas tinjauannya yang luas, berbagai kata baik, dan dorongannya sepanjang
upaya menyusun edisi pertama. Saya juga berutang banyak terima kasih
kepada rekan-rekan di CollegeOpenTextbooks.org, yang menyediakan penyuntingan
naskah tanpa bayaran. Terakhir, terima kasih kepada para pengguna
yang telah memberikan koreksi dan saran perbaikan. Sejumlah masukan
tersebut telah dimasukkan ke dalam edisi ketiga.
